% Instruction file: workspace/instructions/presentation_oct_23.md
% Build from any directory: bash presentations/build_oct_23.sh
\documentclass[aspectratio=169,10pt]{beamer}
\usepackage{graphicx}
\usepackage{booktabs}
\usepackage{array}
\usepackage{tabularx}
\usepackage{amsmath}
% Standard Beamer theme; colors verified against databento.com CSS.
\usetheme{Madrid}
\definecolor{databentopurple}{HTML}{6266E5}
\definecolor{databentocharcoal}{HTML}{323348}
\setbeamercolor{structure}{fg=databentopurple}
\setbeamercolor{normal text}{fg=databentocharcoal,bg=white}
\setbeamercolor{palette primary}{fg=white,bg=databentopurple}
\setbeamercolor{palette secondary}{fg=white,bg=databentopurple!80!black}
\setbeamercolor{palette tertiary}{fg=white,bg=databentocharcoal}
\setbeamercolor{frametitle}{fg=white,bg=databentopurple}
\setbeamercolor{title}{fg=white,bg=databentopurple}
\setbeamertemplate{navigation symbols}{}
\setbeamertemplate{footline}[frame number]
\title{P116: Order Book Dynamics\\Around Major Events}
\author{}
\date{}
% Search both repository-root and presentations-directory working directories.
\graphicspath{{outputs/}{../outputs/}}
\newcommand{\chart}[2]{\begin{center}\includegraphics[width=\linewidth,height=#2,keepaspectratio]{#1}\end{center}}
\newcommand{\takeaway}[1]{{\small\textbf{Observation:} #1}}
\begin{document}

% 1
\begin{frame}[plain]
\centering
\titlepage
\end{frame}

% 2
\begin{frame}{Progress and discussion agenda}
\begin{enumerate}
\item \textbf{Progress to date}
  \begin{itemize}
  \item Explore the platform and establish the data workflow
  \item Select assets and macroeconomic announcements
  \item Collect daily futures data and define contract selection
  \item Compare daily moves on announcement and normal days
  \end{itemize}
\item \textbf{Next research steps}
  \begin{itemize}
  \item Analyse intraday dynamics using market-by-order (MBO) data
  \item Discuss a dashboard for monitoring market conditions
  \end{itemize}
\end{enumerate}
\end{frame}

% 3
\begin{frame}{Asset universe}
\small
Futures offer liquid benchmark exposures, long trading hours and a common exchange data source.
\medskip
{\renewcommand{\arraystretch}{1.25}
\begin{tabularx}{\linewidth}{@{}>{\raggedright\arraybackslash}p{0.23\linewidth}>{\raggedright\arraybackslash}X@{}}
\textcolor{databentopurple}{ZQ: Fed funds} & Current calendar month and next six months; expected average overnight rates and policy-path repricing. \\
\textcolor{databentopurple}{ZT: 2-year Treasury} & Short-rate expectations and demand for safe assets. \\
\textcolor{databentopurple}{ZN: 10-year Treasury} & Longer-run rate expectations and term premium; complements ZT in studying the curve. \\
\textcolor{databentopurple}{GC: Gold} & Responses to real rates, the dollar and demand for safe assets. \\
\textcolor{databentopurple}{6E: EUR/USD} & Dollar-relative currency response; higher prices mean a stronger euro. \\
\textcolor{databentopurple}{ES: S\&P 500} & Equity responses to rates, earnings expectations and growth news. \\
\textcolor{databentopurple}{BTC: Bitcoin} & Compare a newer asset's response with established markets.
\end{tabularx}}
\medskip

All contracts are available through Databento's \texttt{GLBX.MDP3} dataset.
\end{frame}

% 4
\begin{frame}{Daily data and contract selection}
\begin{itemize}
\item Retrieve explicit futures contracts using Databento's \texttt{ohlcv-1d} schema.
\item \textbf{Fed funds:} retain seven calendar-month horizons, T0--T6, on each date.
\item \textbf{Other assets:} retain two nearby contracts and mark one as primary.
\item Switch primary when the challenger accounts for \textbf{at least 65\% of their combined volume for three consecutive trading observations}.
\item If only one contract has a bar, use that contract immediately; if neither has a bar, neither is primary.
\item Compute each open-to-close return within one contract. No return is calculated across a roll-price gap.
\end{itemize}
\end{frame}

% 5
\begin{frame}{Daily analysis before the intraday study}
\small
\textbf{Coverage:} June 2010--October 2026; Bitcoin has a shorter history.
Events: FOMC, CPI, nonfarm payrolls (NFP), and PCE releases.

\textbf{Time window:} bars span midnight to midnight UTC, crossing CME session boundaries. Intraday research will align observations to release timestamps.
\begin{enumerate}
\item \textbf{Same-day correlations:} compare signed open-to-close moves across assets and event samples.
\item \textbf{Lag-one autocorrelation:} correlate today's move with the next observed trading day's move; condition on the first day being an event.
\item \textbf{Range intensity:} divide today's high--low range by its preceding 60-trading-date mean, excluding today.
\end{enumerate}
FF uses PCA of seven rate changes or ranges. Other signed returns are \((C-O)/O\). Normal days exclude all four event categories.
\end{frame}

% 6
\begin{frame}{Fed funds curve PCA}
\chart{ff_curve/ff_curve_pca_all_dates.png}{0.61\textheight}
\small
Covariance PCA: demean each maturity after p5--p95 winsorization; no scaling by maturity standard deviation. Eigenvectors have unit length.

\takeaway{PC1 is a ramp rather than a parallel shift. Further horizons carry more movement variance; current-month averaging may also suppress the front end.}
\end{frame}

% 7
\begin{frame}{Cross-asset correlations: all dates}
\chart{daily/open_close_correlations_all_dates.png}{0.68\textheight}
\takeaway{The rates group is strongly linked: ZT--ZN = 0.80; FF rate-change PC1--ZT = \(-0.78\). All dates include event days.}
\end{frame}

% 8
\begin{frame}{Cross-asset correlations: FOMC dates}
\chart{daily/open_close_correlations_fomc_dates.png}{0.68\textheight}
\takeaway{Several relationships strengthen: BTC--ES 0.34 to 0.63; gold--EUR 0.33 to 0.61; BTC--EUR 0.12 to 0.42. BTC pairs have 72 FOMC observations.}
\end{frame}

% 9
\begin{frame}{Cross-asset correlations: CPI dates}
\chart{daily/open_close_correlations_cpi_dates.png}{0.68\textheight}
\takeaway{FF PC1--ZT = \(-0.87\); gold--EUR = 0.44 and BTC--ES = 0.39. The FOMC pattern is not shared uniformly across releases.}
\end{frame}

% 10
\begin{frame}{Cross-asset correlations: NFP dates}
\chart{daily/open_close_correlations_nfp_dates.png}{0.68\textheight}
\takeaway{Strong links span rates, gold and FX: ZN--gold = 0.58 and gold--EUR = 0.56. BTC--ES = 0.32, close to its full-sample value.}
\end{frame}

% 11
\begin{frame}{Lag-one autocorrelation of daily moves}
\begin{columns}[c]
\begin{column}{0.69\textwidth}
\chart{daily/open_close_lag1_autocorrelations.png}{0.76\textheight}
\end{column}
\begin{column}{0.28\textwidth}
\small
\textbf{Compare:} day \(t\) open-to-close with day \(t+1\) open-to-close.

Select samples by day \(t\); day \(t+1\) may be any event category.

\textbf{Gold stands out:} FOMC \(\rho=0.33\), versus \(-0.04\) on normal days.

Positive values indicate linear continuation relative to sample means, not a guaranteed repeated sign.

FF uses a common full-sample PC1 basis directly, without differencing PC1 again.
\end{column}
\end{columns}
\end{frame}

% 12
\begin{frame}{Daily range intensity: FOMC dates}
\chart{daily/daily_vol_fomc_vs_normal.png}{0.66\textheight}
\small
\(I_t=R_t / \operatorname{mean}(R_{t-60},\ldots,R_{t-1})\).
For FF, \(R_t\) is the uncentered projection of clipped seven-horizon ranges onto range-PC1 loadings.

\takeaway{The largest median increases are in FF and ZT. BTC shows little range expansion despite stronger equity co-movement. This is a range proxy, not a standard-deviation estimate.}
\end{frame}

% 13
\begin{frame}{Daily range intensity: CPI dates}
\chart{daily/daily_vol_cpi_vs_normal.png}{0.70\textheight}
\takeaway{Rate futures show elevated median ranges. Gold, ES and BTC medians are near their trailing means, but their upper quartiles exceed normal-day levels.}
\end{frame}

% 14
\begin{frame}{Daily range intensity: NFP dates}
\chart{daily/daily_vol_nfp_vs_normal.png}{0.70\textheight}
\small
Compare the median and upper quartile with normal days, using the same preceding 60-date baseline. A value of 1 means today's range equals its trailing mean.
\end{frame}

% 15
\begin{frame}{Proposed intraday research}
\begin{columns}[t]
\begin{column}{0.32\textwidth}
\textbf{Liquidity}
\begin{itemize}
\item Bid--ask spread
\item Depth at level \(k\)
\item Order additions and cancellations by level
\end{itemize}
\end{column}
\begin{column}{0.30\textwidth}
\textbf{Order flow}
\begin{itemize}
\item Total traded volume
\item Signed trade flow
\item Arrival and cancellation rates
\end{itemize}
\end{column}
\begin{column}{0.30\textwidth}
\textbf{Prices}
\begin{itemize}
\item Returns around release time
\item Exponentially weighted volatility
\item Price impact of order flow
\end{itemize}
\end{column}
\end{columns}
\vspace{8mm}
These are proposals for discussion; prioritisation will depend on data coverage and reconstruction cost.
\end{frame}

% 16
\begin{frame}{Study design and key metrics}
\begin{enumerate}
\item \textbf{Event-aligned trajectories}
\begin{itemize}
\item Compare pre- and post-release features with normal days at matching times.
\item Measure time to peak or trough and time to stabilisation.
\item Compare recovery speeds across assets and announcement types.
\end{itemize}
\item \textbf{Information flow across markets}
\begin{itemize}
\item Study lead--lag relationships in prices, order flow and liquidity.
\item Does EUR futures order flow precede gold returns or gold order flow?
\item Does liquidity withdrawal in one market precede withdrawal elsewhere?
\item Explore event-specific flow hierarchies: which markets lead, which follow, and whether that ordering differs across announcements.
\end{itemize}
\end{enumerate}
\small
Check clock alignment, influential days and stability across periods before interpreting apparent leads as predictive relationships.
\end{frame}

% 17
\begin{frame}{Dashboard proposal}
\begin{itemize}
\item Monitor spreads, depth, signed order flow and price volatility across the chosen markets.
\item Compare current conditions with historical event and normal-day patterns.
\item Condition reference trajectories on the release type and elapsed time since the announcement.
\item Explore expected time to normalisation, with updates as observations arrive.
\end{itemize}
\vspace{5mm}
\textbf{Scope for discussion:} prioritise the empirical intraday study; build a dashboard if time and live-data access permit.

\textbf{Feedback welcome:} which assets, features and event comparisons should come first?
\end{frame}

% 18
\begin{frame}[plain]
\centering
{\usebeamerfont{title}\color{databentopurple} Thank you}
\end{frame}
\end{document}
